\subsection{11 августа. Озера. Вершина Марс}
\textit{Метеоусловия: Солнечно, ветренно. }

\begin{figure}[h!]
	\centering
	\includegraphics[angle=0, width=0.7\linewidth]{../pics/maps/11_map.jpg}
	\label{fig:03_map}
\end{figure}

В 4:30 был общий подъём и сбор лагеря параллельно с группой Алексея. Сперва собирались с фонариками, так как ещё была луна на небе, после чего стало ясно, без осадков и ветра. В 5:30 вышли чуть раньше дружественной группы. Первый день похода, когда участника n не гоняшит она узнала, что мы находимся в горах. Благодаря этому темп был активнее, чем в прошлые дни
В 6:30 был первый брод, который просто прошли поперёк, дальше, в 7:55, следовал второй брод, сперва искали переправу минут 10, после чего получилось перейти не снимая ботинок и по камням. 

\begin{figure}[h!]
	\centering
	\includegraphics[angle=0, width=0.7\linewidth]{../pics/days/11_1.jpg}
	\label{fig:11_1}
	\caption{ Движение группы по треку}
\end{figure}

\begin{figure}[h!]
	\centering
	\includegraphics[angle=0, width=0.7\linewidth]{../pics/days/11_2.jpg}
	\label{fig:11_2}
	\caption{ Наша группа (руководитель) нашёл череп горного козла и взял его с собой. Звать его Кардинал.}
\end{figure}

Темп был хороший, руководитель решил плавно набирать высоту, именно это и начали делать в 8:45. К 11:15	достигли отметки  4035. 	
\begin{figure}[h!]
	\centering
	\includegraphics[angle=0, width=0.7\linewidth]{../pics/days/11_3.jpg}
	\label{fig:11_3}
	\caption{ Группа на спуске к обеду. Вдали видно вершину Марс и озеро, у которого был обед.}
\end{figure}

В 11:55, высота 3900. 
Обед. Купаемся в озере. Команда в активном и приподнятом настроении устроила соревнование по берпи.  Обедали рядом с группой Алексея Остапива. Разложили и сушили палатки, тенты и мокрые вещи на солнце.

\begin{figure}[h!]
	\centering
	\includegraphics[angle=0, width=0.7\linewidth]{../pics/days/11_4.jpg}
	\label{fig:11_4}
	\caption{Группа идёт вдоль озера к МН8.}
\end{figure}

В 15:25 дошли до места ночевки МН8. Группа из 4-х человек готовится к восхождению на гору Марс, в лагере осталось два участника на рации и с готовностью готовить ужин. Остались по соображениям здоровья, чтобы смочь пройти оставшийся участок похода без ухудшения состояния. Каждый час была связь с лагерем по рации, чтобы удостовериться, что группа идёт хорошо, ничего не случилось и она укладывается в тайминги. У участников на руках был отчёт группы, которая из подъёма в. Марс устроила днёвку и затратила 5 часов на подъём и спуск суммарно. Руководитель прикинул и сказал, что мы в активном темпе точно успеваем спуститься вниз на безопасный и пологий участок до захода солнца.

\begin{figure}[h!]
	\centering
	\includegraphics[angle=0, width=0.7\linewidth]{../pics/days/11_5.jpg}
	\label{fig:11_5}
	\caption{Группа на подъёме.}
\end{figure}
18:37, высота  4300, взошли на вершину Марс.
\begin{figure}[h!]
	\centering
	\includegraphics[angle=0, width=0.7\linewidth]{../pics/days/11_6.jpg}
	\label{fig:11_6}
	\caption{Группа на вершине. Сзади ледовый купол и д.р Арабель-Суу.}
\end{figure}

\begin{figure}[h!]
	\centering
	\includegraphics[angle=0, width=0.7\linewidth]{../pics/days/11_7.jpg}
	\label{fig:11_7}
	\caption{Группа любуется видами и пишет записку для туристов или альпинистов. }
\end{figure}

Пофоткались, нашли свой лагерь с высоты, съели шоколадку, написали записку и вниз, скорость подъёма и спуска выше средней скорости движения группы. Участники хотели спуститься с неприятного участка до захода солнца. У них получилось. За один час участники сбросили 500 метров, то есть, в 20:05 они уже спустились к подножью Марса. В 20:42 вернулись в лагерь, покушали и легли спать, ведь завтра вновь ранний подъем, но позднее, чем если бы на Марс не пошли, чтобы все могли выспаться.
